\frontchapter{摘 要}

三维医学图像分割需要在体数据中同时组织局部解剖细节、跨区域与跨切片的长距离上下文，以及U形网络不同尺度层级中的结构信息。局部边缘、纹理和邻域连续性依赖细粒度空间表征，而较远区域之间的解剖联系需要更大范围的信息交互；与此同时，编码与解码过程中不断变化的空间分辨率和语义层级，又使同一结构在不同尺度下具有不同的表征条件。针对上述问题，本文围绕三维体数据中多类空间信息的协调建模，从空间维度的局部--全局协同和尺度维度的结构信息组织两个相互衔接的层面开展研究。

针对局部空间特征与长距离上下文需要协同建模的问题，本文提出基于Vision-xLSTM局部--全局协同建模的三维医学图像分割方法ViL-UNet。该方法以3D CNN形成局部与层次化三维特征，并将深层特征进一步组织为视觉序列，通过Vision-xLSTM建立较长距离的上下文交互；同时采用多方向序列组织扩展不同空间方向上的信息传播，使局部卷积表征与跨区域、跨切片上下文在统一框架中形成互补。解码阶段逐级恢复空间分辨率，并结合跨层特征完成局部细节、层次语义和深层上下文的融合。实验结果显示，ViL-UNet在Synapse数据集上获得85.12\%的平均DSC和12.49~mm的HD95，在ACDC数据集上获得92.79\%的平均DSC，表明所构建的局部--全局协同框架在两类公开三维医学图像分割任务上具有较好的分割表现。

在空间范围内的局部--全局关系得到组织后，U形网络不同分辨率和语义层级中的结构信息仍需要进一步处理。受控前置比较显示，持续累积跨尺度历史状态后，不同评价维度并未形成方向一致的有利变化，由此将研究重点由增加跨尺度传播转向保持结构信息与当前尺度之间的对应关系。基于这一认识，本文进一步提出基于尺度特异结构证据调制的三维医学图像分割方法。该方法从各尺度当前层次特征中独立构造对应的尺度特异结构证据，使结构信息来源保持明确的层级属性；随后联合当前解码状态与对应结构证据形成通道级自适应调制关系，使结构信息的利用方式与当前重建状态相联系；结构增强分量再以残差形式接入解码表示，在保留基础解码路径的同时参与对应尺度的空间恢复。围绕该尺度信息组织链，本文从总体分割性能、器官级差异、尺度信息组织策略、组成关系及模型复杂度等方面建立分析体系，以考察不同尺度结构信息从形成到利用的完整关系。

本文两项研究围绕同一三维医学图像信息组织问题形成由空间关系到尺度层级关系的递进：前者强调局部结构表征与长距离上下文在空间维度上的功能互补，后者进一步关注不同尺度结构信息与当前解码状态之间的对应和利用关系。通过上述研究，本文形成了从局部三维特征建模、长距离空间上下文组织到尺度层级结构信息利用的连续技术路线，为三维医学图像分割中不同维度信息的协调建模提供了一种统一的研究思路。

\noindent\hei 关键词：\song 三维医学图像分割；Vision-xLSTM；局部--全局协同；尺度结构信息；多尺度建模
